\chapter*{Travaux Pratiques-Programmation Orientée Objet}

\section*{Création d'une classe : pas à pas}
Ouvre l'IDE Spyder et créez un nouveau fichier.\\

\subsection*{Création d'attributs et instanciation d'un objet}
\begin{enumerate}
    \item Créez une classe simple, la classe Carte correspondant à une carte d'un jeu. Par convention, le nom d'une classe commence toujours par une majuscule.
\item Créez une méthode constructeur qui va définir les attributs de la cartes qui seront :
\begin{itemize}
    \item sa valeur 2,3···,10, 11 pour Valet, 12 pour  Dame, 13 pour Roi et 14 pour As ;
    \item sa couleur (rouge ou noir),
    \item son enseigne (carreau, trèfle, coeur ou pique)
\end{itemize}
\item Dans le programme principal, créez une instance de la classe Carte. On veut un 7 de coeur. \textit{Lorsque  l'on  crée  un  objet,  son  constructeur  est  appelé  implicitement  et  l'ordinateur  alloue  de  la mémoire pour l'objet et ses attributs. On peut d'ailleurs obtenir l'adresse mémoire de notre objet créé.}
\item Affichez sont type
\item Affichez la valeur de l'instance, sa couleur et son enseigne.
\item Changez le 7 de coeur en 7 de pique.
\item Créer une carte c2: un 4 de coeur.
\item Modifier la couleur de la carte c2 en la passant à pique.
\item Modifier la valeur de la carte c2 en la passant à 8.
\item Afficher les attributs de c2.
\end{enumerate}

\subsection*{Ajout de méthodes, appel de la méthode sur l'instance d'un objet}

Ouvrez le fichier nommé \textit{cartes.py}\\
\begin{enumerate}
    \item Comparer la classe \textit{Carte} à votre classe \textit{Carte}. Quelles sont les différences ?
    \item Dans ce nouveau fichier écrire une méthode \textit{egalite} qui renvoie True si la carte self et la carte donnée en paramètre sont égaux, False sinon.
    \item compléter la méthode \textit{estSuperieurA} qui renvoie True si la carte self est supérieure à la carte donnée en paramètre, False sinon.
    \item compléter la méthode \textit{estInferieurA} qui renvoie True si la carte self est inférieure à la carte donnée en paramètre, False sinon.
    \item On va maintenant tester nos méthodes : créer trois instances de JeuCartes de trois valeurs différentes et comparez les à l'aide des méthodes créées.
\end{enumerate}

\section*{Création d'une nouvelle classe Jeu de Cartes}

Ouvrez le fichier nommé \textit{jeuCartes.py}. La classe JeuCartes va créer un jeu avec 54 cartes. 

\begin{enumerate}
    \item Compléter le constructeur.
    \item Compléter la méthode melanger en utilisant la fonction shuffle du module random. \url{https://docs.python.org/3/library/random.html}.
    \item Compléter la méthode \textit{distribuerCarte} qui défile la prochaine carte du jeu et l'enlève du jeu.
    \item Dans le programme principal, créer une instance de jeu de cartes
    \item Mélangez ce jeu
    \item Distribuez les cartes et affichez la valeur, l'enseigne et la couleur de la carte distribuée à chaque distribution.
\end{enumerate}